\section{Related Work}

\paragraph{Cognitive Science Based Agent Memory}
Cognitive science characterizes long-term memory as persistent storage, and distinguishing episodic, semantic, and procedural memory~\citep{Atkinson1968HumanMemory,Tulving1972EpisodicSemantic,Squire2004MemorySystems}. In agentic settings, interactions are naturally episodic, from which reusable abstractions can be derived, e.g., factual information such as user preferences can be distilled into semantic memory~\citep{li2025memosmemoryosai}, while action strategies can be abstracted as procedural memory~\citep{wang2024agentworkflowmemory}. These abstractions align with more general knowledge notions: propositional knowledge as factual statements and prescriptive knowledge as generalized goal-directed procedures~\citep{10.1257/aer.103.3.545}.

Building on these insights, recent systems incorporate different memory types to support long-context understanding~\citep{lee2024humaninspiredreadingagentgist}, long-term interactions~\citep{li2025memosmemoryosai}, structured memory organization~\citep{anokhin2024arigraphlearningknowledgegraph,rasmussen2025zeptemporalknowledgegraph}, and experience-driven strategy refinement~\citep{zhao2024expelllmagentsexperiential,wang2024agentworkflowmemory,DBLP:journals/corr/abs-2508-06433}. However, most approaches are task-specific and do not jointly support multiple memory types. In contrast, \plugmem standardizes heterogeneous episodic experience and organizes extracted knowledge into a propositional-prescriptive structure, enabling task-agnostic reuse.

\paragraph{Memory Module Designs}
\begin{table}[ht]
\centering
\small
\caption{Comparison of representative agentic memory systems.
\textbf{M2K} indicates if a system transforms working memories into \textit{reusable, abstract knowledge} that's generalizable across tasks.
\textbf{KaU} means if it adopts knowledge as the memory unit.
\textbf{Mech.} characterizes the design of the memory modules, including
\textbf{Str.} (structured memory representations),
\textbf{Ret.} (refined retrieval strategies),
and \textbf{Rea.} (post-retrieval reasoning or compression).
\textbf{Mem.} denotes the memory type that the system supports,
including \textbf{E}pisodic, \textbf{S}emantic, and \textbf{P}rocedural memory.
\textbf{Agn.} indicates whether the memory design is task-agnostic.
}

\resizebox{\columnwidth}{!}{
\begin{tabular}{lccccccccc}
\toprule
\multirow{2}{*}{\textbf{Meth.}} &
\multirow{2}{*}{\textbf{M2K}} &
\multirow{2}{*}{\textbf{KaU}} &
\multicolumn{3}{c}{\textbf{Mech.}} &
\multicolumn{3}{c}{\textbf{Mem.}} &
\multirow{2}{*}{\textbf{Agn.}} \\
\cmidrule(lr){4-6}\cmidrule(lr){7-9}
&
&
&
\textbf{Str.} &
\textbf{Ret.} &
\textbf{Rea.} &
\textbf{E} &
\textbf{S} &
\textbf{P} &
\\
\midrule
Vanilla Retrieval       & \xmark & \xmark & \xmark & \xmark & \cmark & \cmark & \xmark & \xmark & \cmark \\
Vanilla RAG       & \xmark & \xmark & \xmark & \xmark & \xmark & \cmark & \xmark & \xmark & \cmark \\
GraphRAG          & \cmark & \xmark & \cmark & \cmark & \cmark & \xmark & \cmark & \xmark & \pmark \\
A-Mem             & \cmark & \cmark & \cmark & \xmark & \xmark & \cmark & \cmark & \xmark & \cmark \\
\midrule
Zep               & \xmark & \xmark & \cmark & \cmark & \xmark & \cmark & \cmark & \xmark & \xmark \\
MemoryOS          & \xmark & \xmark & \xmark & \cmark & \xmark & \cmark & \cmark & \xmark & \xmark \\
HippoRAG2         & \cmark & \xmark & \cmark & \cmark & \cmark & \xmark & \cmark & \xmark & \xmark \\
AWM               & \cmark & \cmark & \xmark & \xmark & \xmark & \xmark & \xmark & \cmark & \xmark \\
ReasoningBank     & \cmark & \cmark & \xmark & \xmark & \xmark & \cmark & \xmark & \cmark & \xmark \\
\specialrule{0.6pt}{0pt}{2pt}
\rowcolor{ShadowBlue}
\textbf{\plugmem} & \cmark & \cmark & \cmark & \cmark & \cmark & \cmark & \cmark & \cmark & \cmark \\
\bottomrule
\end{tabular}
}
\vspace{4pt}
\label{tab:related_works}
\end{table}

Many agent memory systems adopt retrieval-based external memory, storing past interactions and retrieving them at inference time. Early task-agnostic methods rely on flat retrieval over unstructured episodic memory, such as vanilla retrieval and RAG~\citep{lewis2021retrievalaugmentedgenerationknowledgeintensivenlp,Wang_2024}, resulting in redundant and episode-specific memory reuse. Later work introduces structure over episodic memory, including hierarchical and graph-based retrieval (e.g., GraphRAG) to support multi-hop reasoning~\citep{edge2025localglobalgraphrag,sarthi2024raptorrecursiveabstractiveprocessing,gutiérrez2025ragmemorynonparametriccontinual,Wang_2025}. While improving access efficiency, these methods largely preserve episodic traces as the primary memory unit.
In parallel, task-specific systems explicitly transform experience into higher-level representations, such as temporal knowledge graphs or workflow memories~\citep{gutiérrez2025ragmemorynonparametriccontinual,wang2024agentworkflowmemory,ouyang2025reasoningbankscalingagentselfevolving}. Although effective within fixed tasks, their abstractions are tightly coupled to task assumptions and limit transferability.

Overall, as shown in Table \ref{tab:related_works}, prior work highlights a distinction between improving retrieval over episodic memory and transforming experience into reusable knowledge, with only the latter supporting cross-task generalization. \plugmem adopts this principle by organizing memory around propositional and prescriptive knowledge units, enabling more effective retrieval of decision-relevant information, while retaining episodic traces as verifiable evidence.

For a more detailed discussion of related work, including the full versions of the above sections
and additional discussion on memory benchmarks, see Appendix~\ref{app:related_work}.
